\subsection{Duality and Weak Convergence}\label{sec:duality-weak-convergence}

We have already encountered many of the ideas of duality (studying a function by studying how linear functionals act upon it) in the theory of distributions. These ideas are very powerful in the study of Banach spaces and Hilbert spaces as well.

\subsubsection{Bounded Linear Mappings}\label{subsec:bounded-linear-mappings}

\begin{Definition}
\label{def:bounded-linear-mapping}
Let $X,Y$ be normed vector spaces. A linear mapping $L:X\to Y$ is called \textbf{bounded} if there is a constant $C$ such that $\|Lx\|\leq C\|x\|$ for every $x\in X$.
\end{Definition}

For linear mappings, the ideas of continuity and boundedness are essentially the same.

\begin{Definition}
\label{def:operator-sequential-continuity}
We say that an operator $L:X\to Y$ is \textbf{continuous} at a point $x\in X$ if, whenever $x_n\in X$ is a sequence such that $x_n\to x$, we have $L(x_n)\to L(x)$.
\end{Definition}

\begin{Remark}
\label{rem:sequential-versus-topological-continuity}
Definition~\ref{def:operator-sequential-continuity} is often called the definition of \emph{sequential continuity} to distinguish it from the topological version of the definition.
\end{Remark}

\begin{Lemma}
\label{lem:linear-mapping-bounded-continuity-equivalence}
Let $L:X\to Y$ be a linear mapping.\begin{enumerate}
\item If $L$ is continuous at the origin, it is continuous at every $x\in X$.
\item $L$ is continuous if and only if it is bounded.
\end{enumerate}
\end{Lemma}

\begin{Proof}
To prove part 1 we assume that $L$ is continuous at the origin. Note that the linearity of the operator implies $L(0)=0$. Then for any $x\in X$ and any sequence $x_n\in X$ with $\lim_{n\to\infty}x_n=x$ we have
\begin{equation}\label{eq:linear-map-continuity-at-origin-translation}
\lim_{n\to\infty}L(x_n)-L(x)=\lim_{n\to\infty}L(x_n-x)=L(0)=0,
\end{equation}
so $L$ is continuous at $x$.

To prove part 2 we note that the definition of boundedness implies continuity at zero. Thus, by part 1, boundedness implies continuity. On the other hand, suppose $L$ is not bounded. Then for every $n$ there exists $y_n$ such that $\|Ly_n\|>n^2\|y_n\|$. Note that $x_n:=y_n(n\|y_n\|)^{-1}\to0$ but $\|Lx_n\|>n$. It follows that $L$ is not continuous at the origin. Thus, continuity implies boundedness.
\end{Proof}

It is natural to consider the set of all bounded linear mappings. It follows immediately from the definition that this set forms a vector space. Also, if we take the smallest possible constant in Definition~\ref{def:bounded-linear-mapping}, then this quantity gives us a measure of the ``size'' of a linear mapping. This motivates the following definition.

\begin{Definition}
\label{def:bounded-operator-space-norm}
By $\mathcal{L}(X,Y)$ we denote the set of all bounded linear mappings from $X$ to $Y$. If $X=Y$, we also write $\mathcal{L}(X)$ for $\mathcal{L}(X,X)$. Moreover, if $L\in\mathcal{L}(X,Y)$, we set
\begin{equation}\label{eq:bounded-linear-map-operator-norm}
\|L\|:=\sup_{\|x\|=1}\dfrac{\|Lx\|}{\|x\|}=\sup_{\|x\|\ne0}\dfrac{\|Lx\|}{\|x\|}.
\end{equation}

\end{Definition}

\begin{Theorem}
\label{thm:bounded-operator-space-completeness}
Let $X,Y$ be Banach spaces. Then $\mathcal{L}(X,Y)$, with the norm defined by \eqref{eq:bounded-linear-map-operator-norm}, is also a Banach space.
\end{Theorem}

The proof is straightforward and is left as an exercise (Problem~\ref{ex:bounded-operator-space-completeness}). Linear mappings, also called linear operators, will be studied more extensively in the next chapter. In this section, we are interested in the special case of linear mappings from a Banach space to its scalar field.

\begin{Definition}
\label{def:linear-functional-dual-space}
Let $X$ be a real (complex) Banach space. Then a \textbf{linear functional} on $X$ is a bounded linear mapping from $X$ to $\mathbb{R}$ ($\mathbb{C}$). The space of all linear functionals on $X$ is called the \textbf{dual space} of $X$ and is denoted by $X^*$.
\end{Definition}
